\phantomsection\label{sec:research-plan}
\section*{研究计划}
\subsection*{已有研究方向：时空数据挖掘}
我已经建立的研究方向是时空数据挖掘，即从带有空间信息且随时间变化的大规模数据中提取并利用规律。
这类数据通常来自交通场景。
表示学习是贯穿我各项工作的主线，其目标是学习不依赖具体任务的数据表示，从而避免为每项任务单独构建模型。
表示学习是通向时空数据挖掘基础模型的一条有前景的路径，这类基础模型能够泛化到不同的数据来源和任务，我的 UniTE、UVTM 和 TransferTraj 等工作已经印证了这一点。

\subsection*{新兴研究方向：人工智能材料科学}
在已有研究的基础上，我正在开拓将数据挖掘用于材料科学的跨学科方向，目标是推动人工智能材料科学的发展。
玻璃等材料本身可以视为具有动态结构的空间数据，挖掘这类数据需要运用我为时空数据开发的同一套表示学习方法，因此这一方向是我已有研究的自然延伸。
我已经通过 Villum 资助的博士后项目“基于扩散概率模型的材料逆向设计”开展这一方向的研究，并与 Morten M. Smedskjaer 研究组合作。
该研究组是无序材料科学领域的领先团队之一。

\vspace{\entryskip}
这一方向的前景在于，人工智能可以在广阔的设计空间中发现具有定制性能的新材料，而传统试错方法探索这一空间的速度太慢、成本太高。
我在这一方向已经提出用于非晶材料逆向设计的扩散模型 AMDEN。
相关论文于 2026 年发表于 Advanced Materials，我为共同第一作者。
我计划深化并拓展这一方向，成长为人工智能材料科学领域一名专注于无序材料的领先研究人员。
无序材料这一领域具有突出的研究和应用潜力，但尚未得到充分探索。

\subsection*{社会影响与应用前景}
玻璃等无序材料支撑着储能、能源生产和先进光学等众多技术。
人工智能可以替代大量耗费资源的实验室试错实验，加快发现具有定制性能的新型无序材料。
更快地发现这些材料有助于实现欧盟和联合国的绿色转型目标，包括负责任消费和生产、气候行动等联合国可持续发展目标。
与我合作的无序材料研究组能够在实验室中制备并验证有潜力的候选材料，使成功的设计有机会走向实际应用。

\vspace{\entryskip}
这一方向连接计算机科学与材料科学，也将形成持久的跨学科研究能力，并促进两个领域之间的合作，使双方长期受益。
